\section{Introduction}\label{sec:intro}
A connected triple cover of a genus-two curve, simply branched at two
values, has genus five. Suppose the curve is embedded in a symplectic
six-manifold and we vary its normal Cauchy--Riemann operator while
keeping the curve fixed. A deformation of the cover away from the
curve becomes possible when its pullback normal operator acquires a
kernel. We show that passage through this wall can change the local
curve count both in the class of the triple cover and in degree seven. Every
connected unramified cover through degree seven
can have the same determinant sign on both sides.

To specify the local count, let \(C\) have genus \(g\ge2\) and let
\(N\to C\) be a complex rank-two bundle of degree \(2g-2\), as for the
normal bundle of a curve in a threefold of vanishing first Chern class.
A real Cauchy--Riemann operator \(D\) on \(N\) defines an almost
complex structure on its total space. Over a holomorphic map \(u\),
a section \(\xi\) of \(u^*N\) gives a holomorphic graph precisely when
\(u^*D\xi=0\).
Call \(D\) rigid through degree \(m\) if every connected smooth finite
cover of degree at most \(m\) has zero pullback kernel, and super-rigid
if this holds in all degrees. On an unramified cover the normal
operator has index zero. At a rigid endpoint its determinant has
a sign, taken with the orientation continued from complex-linear operators.

The contribution of an isolated curve in class \([C]\) has an integral
decomposition into the elementary clusters of Ionel and Parker
\cite{ip,diw,WueProduct}. We write \(n_{d,h}\) for the signed multiplicity
of the elementary factor of genus \(h\) in class \(d[C]\) in the
isolating neighborhood. Its normalization is the local series of a
genus-\(h\) curve with opposite normal weights. Bai and Swaminathan coorient the wall of smooth double covers and
describe the bifurcation of a smooth cubic cover
\cite[Theorem~4.12(ii)]{bai}. We realize the cubic crossing on a
prescribed cover and compute the higher-degree change due to its nodal
joins. Doan and Walpuski study chambered invariants of real
Cauchy--Riemann operators \cite{dw}. Along paths of the realized unramified crossings, the local series on
intermediate covers determine the change on the original curve
\cite[Theorems~4.1 and~7.1]{WueUnramified}. We ask what the unramified determinant signs can miss
at a branched crossing.

At an ordinary crossing the pullback kernel is a real line with
embedded graph. Differentiating the pullback equation on a generator
in the branch values and the path parameter identifies those
parameters with the cokernel. We write \(\Delta\) for the value after
a crossing minus the value before it.

\begin{theorem}\label{thm:invisible}
On any genus-two curve there is an ordinary crossing on a prescribed connected \(S_3\) triple cover with two distinct simple branch values, with super-rigid endpoints arbitrarily close on opposite sides. All determinant signs on the connected unramified covers of degree at most seven agree at these endpoints. The local elementary changes satisfy
\[
       |\Delta n_{3,5}|=2,\qquad |\Delta n_{7,13}|=1 .
\]
\end{theorem}

The cover of degree three has trivial deck group. A displacement produces one embedded graph on each side of the
crossing. Their fixed-parameter normal signs are opposite. Subtracting
the incoming count therefore gives two contributions of the same sign. In degree seven a normal height on the
sextic Galois closure vanishes above each branch value, where its
graph can join the zero section.
The proof of Theorem~\ref{thm:invisible} is completed in
\S\ref{p14-sec:chambers}.

\subsection{The wall in degree seven}
For a rigid operator, matching across nodes also gives zero kernel
on connected stable maps in the same degree. Its oriented cokernels
form a real bundle \(V_D\) on the compact moduli stack
\(\overline{\mathcal M}_k(C,d)\) of unmarked connected degree-\(d\)
stable maps from curves of arithmetic genus \(k\). Its real rank
is twice the virtual complex dimension \(2k-2-2d(g-1)\). Set
\[
 L_{d,k}(D)=\int_{[\overline{\mathcal M}_k(C,d)]^{\rm vir}}e(V_D),
 \qquad L_d(D;v)=\sum_{k\ge0}L_{d,k}(D)v^{2k-2},
 \qquad x=(2\sin(v/2))^2.
\]
The normal obstruction comparison identifies this Euler integral with
the local threefold invariant \cite[Theorem~1.2]{zinger}. An elementary
factor of genus \(h\) contributes \(x^{h-1}\) in its first degree
\cite[Theorem~2]{pandhodge}. This distinction between the domain genus
\(k\) and the elementary genus \(h\) will matter whenever a stable map
has contracted components. Write \(H_d=1+d(g-1)\) for the genus of
an unramified degree-\(d\) cover.

Fix a connected simply branched cubic \(f:Y\to C\) with monodromy
\(S_3\). Write \(E\) for its set of branch values, \(b=|E|=2\ell>0\),
and \(h_j=1+7(g-1)+3\ell+j\). The sextic closure has genus
\(1+6(g-1)+3\ell\). Attaching it to \(C\) at \(j\) distinct points
gives precisely the arithmetic genus \(h_j\).

\begin{theorem}\label{p14-thm:wall}
At a generic transverse cubic crossing there is one punctured interval
on which every pullback to a connected stable map of degree seven is
injective, in every source genus. Choose super-rigid endpoints on
opposite sides in this interval and write \(\Delta\) for the value on
the positive side minus the value on the negative side. There are real
weights \(w_i\) indexed by \(E\), with all nonempty subsums nonzero, and
one orientation sign \(\kappa\) such that, with
\(a_I=\kappa(-1)^{|I|}\sgn(\sum_{i\in I} w_i)\),
\begin{align}
 \sum_{k\geq0}\Delta L_{7,k}\,v^{2k-2}
     &=\sum_{\varnothing\ne I\subseteq E}a_I x^{h_{|I|}-1},
            \label{p14-eq:main-series}\\
 \Delta n_{7,h_j}&=\sum_{|I|=j}a_I .\label{p14-eq:main-exponents}
\end{align}
All other degree-seven elementary multiplicities are constant across
the crossing. For fixed \(f\) and \(N\), every chamber of the real hyperplanes
where a nonempty subsum vanishes occurs in this way. The common sign can equivalently be absorbed
by reversing all the weights.
\end{theorem}

The interval assertion concerns degree exactly seven. For the endpoint
factorizations we also need rigidity in other degrees, which is a
residual condition on operators \cite{wendl,dwbn}. On a generic path
through the prescribed crossing one may ensure super-rigidity at a
countable dense set of noncentral times on both sides: evaluation at
each of these times is a submersion on paths with fixed first jet at
zero, by variations supported away from zero. Intersect the corresponding
countably many residual conditions. Rigidity in all degrees on the
whole punctured interval is unnecessary.


\subsection{The sign at the nodes}
The three sheet values of a cubic kernel section have sum zero. On
the sextic Galois closure their conjugates span the real standard
representation of \(S_3\). Fix one transposition \(\tau\). Over each branch value the three
ramification points of the sextic have the three different inertia
transpositions. Exactly one is fixed by \(\tau\). At each of these
chosen points a height in the negative eigenspace of \(\tau\) vanishes:
its value equals its own negative. A nonzero vector on this line has
trivial stabilizer, so its graph requires all six sheets. In
degree seven only one sheet remains to carry the zero height.
Selecting a nonempty set of joins connects the two components.

At a chosen node there are two complex variables: one moves the
attaching point along the graph and the other smooths the node.
Modulo the tangent of the graph, the leading matching term is
antilinear in the smoothing variable. It contributes one minus sign
per node. The other line of the standard representation, fixed by the
transposition, supplies an even height on the entire sextic component.
Every matching equation sees this same real parameter. After the
complex variables are eliminated there is one real equation left.
Pairing the next term at each node with the cubic adjoint gives a real
weight \(w_i\). The coefficient of the common height is, up to common
sign, \(\sum_{i\in I}w_i\). In particular one cannot recover the sign
of a two-node join from the individual node signs. For two branch
values on a genus-two target the four possible pairs
\((\Delta n_{7,12},\Delta n_{7,13})\) are
\[
             (-2,1),\qquad (0,1),\qquad (0,-1),\qquad (2,-1).
\]
A displacement replaces each isolated join by regular graphs.
Degree-one maps to these graphs account for contracted tails and bridges
and give the full factor \(x^{h_j-1}\).

\subsection{Compensation on a compact target}
When the normal model is realized on a closed threefold, deformation
invariance constrains the curves outside its isolating tube. At each
generic endpoint perturbation the bounded-energy cycle space splits
into parts supported in the tube and in its complement. Call the
elementary multiplicities of the latter part complementary.
\begin{theorem}\label{p14-thm:torus}
On the standard symplectic six-torus, fix a genus-\(g\) symplectic
curve of primitive class \(A\) as in \S\ref{p14-sec:compact}. For each
positive even \(b\) the cubics in Theorem~\ref{p14-thm:wall} with \(b\)
branch values can be realized on this curve. The changes in the
complementary elementary multiplicities in class \(7A\) are the
negatives of \eqref{p14-eq:main-exponents}. The sharp maximum of the sum
of their absolute values over all genera is \(2^b-1\).
\end{theorem}


The exterior change is determined in each integral class and genus
even when the curves realizing it differ at the two endpoints. With
all branch weights of the same sign, all subsets of a given size
contribute with the same sign, which gives the sharp bound.

The smooth-cover calculation and the first nodal comparison, where
a double cover joins the original curve, occupy
\S\S\ref{sec:chamber-statement}--\ref{sec:cubic}. The application to two spin endpoints in \S\ref{p10-sec:examples}
uses the absolute series of \cite[Theorem~1.1]{WueSpin}. For the
degree-seven theorem we classify the central kernels in
\S\ref{sec:heptic}, compute the local degree at the joins in
\S\ref{p14-sec:joins}, and count contracted components in
\S\ref{p14-sec:virtual}. Sections~\ref{sec:global}
and~\ref{p14-sec:compact} give the smooth degree-seven covers near a
change of sign and the compact compensation. Appendix~\ref{sec:normalization}
compares two normalizations at a double wall, using the compact
example of \cite{WueCounter}. Appendix~\ref{sec:quartic} treats the
first possible genera in higher degrees. Simultaneous crossings of
several branched covers remain open.
